\subsection{Problem Setup}

Let \(G=(\mathcal{B},\mathcal{V},E)\) be a bipartite graph with disjoint node sets
\(\mathcal{B}=\{b_1,\ldots,b_n\}\) and \(\mathcal{V}=\{v_1,\ldots,v_m\}\). Two networks are defined over
these node sets and built from different evidence: a binary network \(Y\in\{0,1\}^{n\times m}\) and a
weighted network \(W\in[0,1]^{n\times m}\). Each pair \((b_i,v_j)\) carries a feature vector
\(z_{ij}\in\mathbb{R}^{d}\).

The first model treats the two networks separately and conditions on the pair features. For each
network the conditional expectation given the features are,
\begin{equation}\label{eq:condmodels}
 \mu_{Y}(z)=\mathbb{E}\!\left[Y_{ij}\mid z_{ij}=z\right],
 \qquad
 \mu_{W}(z)=\mathbb{E}\!\left[W_{ij}\mid z_{ij}=z\right].
\end{equation}
We estimate each conditional expectation within a model class and test it on pairs held out from
fitting.

The second model is a joint prediction model where both networks are predicted from one representation of the same features,
\begin{equation}\label{eq:jointmodel}
 z_{ij}\;\longmapsto\;H_{ij}=h(z_{ij})\;\longmapsto\;\bigl(\widehat Y_{ij},\,\widehat W_{ij}\bigr),
 \qquad H_{ij}\in\mathbb{R}^{r},\quad r<d.
\end{equation}
The question is whether such a shared representation exists: whether the joint model matches the
separate ones in accuracy, and whether the representation it fits is common to both networks rather
than the one a single network induces on its own.

\subsection{General Analysis Pipeline}
\label{sec:pipeline}

The two node sets of the bipartite graph carry their own features:
\(X^{b}\in\mathbb{R}^{n\times d_b}\) holds the feature vectors of the nodes in \(\mathcal{B}\), one per
row, and \(X^{v}\in\mathbb{R}^{m\times d_v}\) those of the nodes in \(\mathcal{V}\). Write \(x^{b}_{i}\)
and \(x^{v}_{j}\) for the rows belonging to \(b_i\) and \(v_j\).

A transformation \(\varphi:\mathbb{R}^{d_b}\times\mathbb{R}^{d_v}\to\mathbb{R}^{d}\) maps the two node
vectors of a pair to a single pair vector,
\begin{equation}\label{eq:pairfeature}
 z_{ij}=\varphi\!\left(x^{b}_{i},x^{v}_{j}\right)\in\mathbb{R}^{d},
\end{equation}

Every pair \((i,j)\) in \(\mathcal{P}=\{1,\ldots,n\}\times\{1,\ldots,m\}\) therefore has a feature vector
\(z_{ij}\) and the two edge values \(Y_{ij}\) and \(W_{ij}\). Stacking over the \(nm\) pairs gives a design
matrix \(Z\in\mathbb{R}^{nm\times d}\) with one row per pair, together with the two networks vectorised into vectors of length
\(nm\).

The problem is then to find mappings from a row of \(Z\) to each edge value, and one mapping that produces both
from a single representation of that row. Figure~\ref{fig:setup} shows the construction.

\begin{figure}[htbp]
  \centering
  \begin{tikzpicture}[
      x = 1cm, y = 1cm, font = \small,
      mat/.style  = {draw = black!55, fill = black!4, line width = .4pt},
      hi/.style   = {draw = black!55, fill = black!30, line width = .4pt},
      gl/.style   = {black!22, line width = .3pt},
      arr/.style  = {-{Stealth[length = 2.2mm]}, line width = .7pt},
      dim/.style  = {font = \scriptsize, text = black!55, inner sep = 1.5pt},
      note/.style = {font = \scriptsize, inner sep = 2pt},
      stage/.style= {font = \footnotesize\itshape, anchor = west, text = black!75}
    ]

  % ---------- (a) node features ----------------------------------------
  \node[stage] at (0,3.55) {(a) node features};

  \draw[mat] (1.5,2.50) rectangle (4.0,3.05);
  \draw[hi]  (1.5,2.68) rectangle (4.0,2.80);
    \node[dim, left]   at (1.40,2.74) {$x^{b}_{i}$};
    \node[dim, above]  at (2.75,3.08) {$d_b$};
    \node[note, below] at (2.75,2.44) {$X^{b}\in\mathbb{R}^{n\times d_b}$};

  \draw[mat] (1.5,0.70) rectangle (4.4,1.55);
  \draw[hi]  (1.5,1.00) rectangle (4.4,1.12);
    \node[dim, left]   at (1.40,1.06) {$x^{v}_{j}$};
    \node[dim, above]  at (2.95,1.58) {$d_v$};
    \node[note, below] at (2.95,0.64) {$X^{v}\in\mathbb{R}^{m\times d_v}$};

  % ---------- (b) the transformation -----------------------------------
  \node[stage] at (6.5,3.55) {(b) one pair};

  \draw[arr] (4.15,2.74) -- (6.60,2.20);
  \draw[arr] (4.55,1.06) -- (6.60,1.70);
  \draw[mat, draw = black!55] (6.95,1.55) rectangle (8.15,2.35);
  \node[note] at (7.55,1.95) {$\varphi$};

  \draw[arr] (8.30,1.95) -- (9.20,1.95);
  \draw[hi] (9.40,1.89) rectangle (12.5,2.01);
    \node[note, right] at (12.6,1.95) {$z_{ij}\in\mathbb{R}^{d}$};
  \node[note] at (10.95,1.30) {$z_{ij}=\varphi\!\left(x^{b}_{i},x^{v}_{j}\right)$};

  % ---------- (c) design matrix ----------------------------------------
  \node[stage] at (0,-0.35) {(c) one row per pair};

  \draw[mat] (1.6,-3.30) rectangle (5.5,-0.90);
  \foreach \y in {-1.30,-1.70,...,-2.95} \draw[gl] (1.6,\y) -- (5.5,\y);
  \foreach \x in {2.05,2.50,...,5.20} \draw[gl] (\x,-3.30) -- (\x,-0.90);
  \draw[hi] (1.6,-2.50) rectangle (5.5,-2.10);
    \node[dim, above] at (3.55,-0.87) {$d$ columns};
    \node[dim, left]  at (1.50,-1.10) {$(1,1)$};
    \node[dim, left]  at (1.50,-1.50) {$(1,2)$};
    \node[dim, left]  at (1.50,-1.90) {$\vdots$};
    \node[dim, left]  at (1.50,-2.30) {$(i,j)$};
    \node[dim, left]  at (1.50,-2.70) {$\vdots$};
    \node[dim, left]  at (1.50,-3.10) {$(n,m)$};
  \node[note, align = center] at (3.55,-3.72)
        {$Z\in\mathbb{R}^{nm\times d}$\\[1pt] row $(i,j)$ is $z_{ij}$};

  \draw[mat] (5.75,-3.30) rectangle (6.10,-0.90);
    \node[dim, above] at (5.92,-0.87) {$y$};
  \draw[hi] (5.75,-2.50) rectangle (6.10,-2.10);
  \draw[mat] (6.35,-3.30) rectangle (6.70,-0.90);
    \node[dim, above] at (6.52,-0.87) {$w$};
  \draw[hi] (6.35,-2.50) rectangle (6.70,-2.10);
  \node[note] at (6.22,-3.62) {edge values};

  % ---------- (d) models ------------------------------------------------
  \node[stage] at (7.35,-0.35) {(d) models};

  \draw[arr] (6.90,-2.30) -- (7.70,-2.30);
  \node[mat, draw = black!55, align = left, inner sep = 5pt, font = \scriptsize]
        at (9.15,-2.30)
        {penalised linear\\ boosted trees\\ neural network\\[2pt] joint model};
  \draw[arr] (10.60,-2.30) -- (11.40,-2.30);

  \draw[mat] (11.60,-3.30) rectangle (11.95,-0.90);
  \draw[hi]  (11.60,-2.50) rectangle (11.95,-2.10);
    \node[note, below, align = center] at (11.77,-3.36) {fitted\\ values};

  \end{tikzpicture}
  \caption[Construction of the prediction problem]{Construction of the prediction problem. Each node
  carries a feature vector. For a pair \((b_i,v_j)\) the transformation \(\varphi\) combines the two node
  vectors into one pair vector \(z_{ij}\). Every pair contributes one row of the design matrix \(Z\), whose
  \(d\) columns are the coordinates of \(\varphi\); the two edge values are aligned with those rows. Each
  model maps a row to one fitted value.}
  \label{fig:setup}
\end{figure}
\FloatBarrier


\subsection{Models}
\label{sec:models}

\subsubsection{Baseline Models}

Every model maps one row \(z_{ij}\in\mathbb{R}^{d}\) to one value, and is fitted on a training index set
\(\mathcal{T}\subset\mathcal{P}\) and evaluated on \(\mathcal{P}\setminus\mathcal{T}\). For a binary edge value \(T_{ij}\in\{0,1\}\) the output is a fitted probability
\(\widehat p_{ij}\in(0,1)\); since the positive class is rare, each observation is weighted inversely to
the frequency of its class in \(\mathcal{T}\), giving weights \(\omega_{ij}>0\) and the weighted
cross-entropy
\begin{equation}\label{eq:bce}
 \ell(\theta)=-\sum_{(i,j)\in\mathcal{T}}\omega_{ij}
 \Bigl[T_{ij}\log\widehat p_{ij}(\theta)+\bigl(1-T_{ij}\bigr)\log\bigl(1-\widehat p_{ij}(\theta)\bigr)\Bigr].
\end{equation}
For a continuous edge value \(\widetilde W_{ij}\in\mathbb{R}\) the output \(\widehat{\widetilde W}_{ij}\) is
on the same scale and the criterion is the sum of squared errors
\begin{equation}\label{eq:sse}
 \mathrm{SSE}(\theta)=\sum_{(i,j)\in\mathcal{T}}
 \bigl(\widetilde W_{ij}-\widehat{\widetilde W}_{ij}(\theta)\bigr)^{2}.
\end{equation}
We use three model families, each in a classification and a regression variant.

\begin{itemize}

\item[--] \textbf{Penalised linear model.} For classification, \(L_{1}\)-penalised logistic regression
\citep{tibshirani1996}: with \(\beta\in\mathbb{R}^{d}\), \(\beta_{0}\in\mathbb{R}\) and
\(\sigma(u)=(1+e^{-u})^{-1}\), the output is \(\widehat p_{ij}=\sigma(\beta_{0}+z_{ij}^{\top}\beta)\) and
\begin{equation}\label{eq:lasso}
 \bigl(\widehat\beta_{0},\widehat\beta\bigr)
 =\operatorname*{arg\,min}_{\beta_{0},\beta}
 \Bigl\{\ell(\beta_{0},\beta)+\lambda\lVert\beta\rVert_{1}\Bigr\},
\end{equation}
and we select \(\lambda>0\) by an inner cross-validation on \(\mathcal{T}\) under the one-standard-error
rule. The \(L_{1}\) penalty sets most of the \(d\) coefficients to exactly zero. For regression, ridge 
regression \citep{hoerl1970}: the output is \(\widehat{\widetilde W}_{ij}=\beta_{0}+z_{ij}^{\top}\beta\) and
\begin{equation}\label{eq:ridge}
 \bigl(\widehat\beta_{0},\widehat\beta\bigr)
 =\operatorname*{arg\,min}_{\beta_{0},\beta}
 \Bigl\{\mathrm{SSE}(\beta_{0},\beta)+\alpha\lVert\beta\rVert_{2}^{2}\Bigr\},
 \qquad \alpha>0.
\end{equation}

\item[--] \textbf{Gradient-boosted trees.} The nonlinear tree model \citep{chen2016}, competitive on tabular feature
spaces \citep{grinsztajn2022}. From a constant \(f_{0}\), regression trees
\(\mathcal{T}_{t}:\mathbb{R}^{d}\to\mathbb{R}\) are fitted successively to the gradient of the criterion,
giving
\begin{equation}\label{eq:xgb}
 f(z)=f_{0}+\eta\sum_{t=1}^{n_{\mathrm{tree}}}\mathcal{T}_{t}(z),
\end{equation}
with \(n_{\mathrm{tree}}\) trees and learning rate \(\eta\in(0,1]\). For classification the trees are fitted to the gradient 
of the cross-entropy and the output is
\(\widehat p_{ij}=\sigma(f(z_{ij}))\) whereas for regression to the gradient of the squared-error criterion
and the output is \(\widehat{\widetilde W}_{ij}=f(z_{ij})\).

\item[--] \textbf{Neural network.} A single hidden layer of width \(r\) \citep{Bishop2006}. With
\(\Psi\in\mathbb{R}^{r\times d}\) and \(\psi\in\mathbb{R}^{r}\),
\begin{equation}\label{eq:mlp}
 H_{ij}=\mathrm{Drop}_{\rho}\!\Bigl(\mathrm{ReLU}\bigl(\Psi z_{ij}+\psi\bigr)\Bigr)\in\mathbb{R}^{r},
\end{equation}
where \(\mathrm{ReLU}(u)=\max(u,0)\) acts coordinatewise and \(\mathrm{Drop}_{\rho}\) zeroes each
coordinate independently with probability \(\rho\) during fitting and is the identity at prediction time.
A linear head with \(\phi\in\mathbb{R}^{r}\) and \(\phi_{0}\in\mathbb{R}\) follows: for classification
\(\widehat p_{ij}=\sigma(\phi^{\top}H_{ij}+\phi_{0})\), minimising the cross-entropy; for regression
\(\widehat{\widetilde W}_{ij}=\phi^{\top}H_{ij}+\phi_{0}\), minimising the squared-error criterion.

\end{itemize}





\subsubsection{Latent-Space Model}

The joint model uses the same hidden layer as the neural network above, but lets one representation feed
two heads. With \(\Psi\in\mathbb{R}^{r\times d}\) and \(\psi\in\mathbb{R}^{r}\),
\begin{equation}\label{eq:trunk}
 H_{ij}=\mathrm{Drop}_{\rho}\!\Bigl(\mathrm{ReLU}\bigl(\Psi z_{ij}+\psi\bigr)\Bigr)\in\mathbb{R}^{r},
\end{equation}
and since \(r<d\) this layer is a bottleneck: the \(d\) coordinates of \(z_{ij}\) are compressed to \(r\)
before either edge value is predicted. Two linear heads act on the same \(H_{ij}\),
\begin{equation}\label{eq:heads}
 \widehat p^{\,Y}_{ij}=\sigma\!\left(\phi_{Y}^{\top}H_{ij}+\phi_{0Y}\right),
 \qquad
 \widehat p^{\,W}_{ij}=\sigma\!\left(\phi_{W}^{\top}H_{ij}+\phi_{0W}\right),
\end{equation}
with \(\phi_{Y},\phi_{W}\in\mathbb{R}^{r}\) and \(\phi_{0Y},\phi_{0W}\in\mathbb{R}\).

Both edge values are binary in this model, so each head is scored by the weighted binary cross-entropy.
Writing \(T^{W}_{ij}\in\{0,1\}\) for the binarised edge weight and \(\theta\) for all parameters,
\begin{equation}\label{eq:bceY}
 \ell_{Y}(\theta)=-\sum_{(i,j)\in\mathcal{T}}\omega^{Y}_{ij}
 \Bigl[\,Y_{ij}\log\widehat p^{\,Y}_{ij}
 +\bigl(1-Y_{ij}\bigr)\log\bigl(1-\widehat p^{\,Y}_{ij}\bigr)\Bigr],
\end{equation}
\begin{equation}\label{eq:bceW}
 \ell_{W}(\theta)=-\sum_{(i,j)\in\mathcal{T}}\omega^{W}_{ij}
 \Bigl[\,T^{W}_{ij}\log\widehat p^{\,W}_{ij}
 +\bigl(1-T^{W}_{ij}\bigr)\log\bigl(1-\widehat p^{\,W}_{ij}\bigr)\Bigr],
\end{equation}
where the weights \(\omega^{Y}_{ij}\) and \(\omega^{W}_{ij}\) are formed separately for each network, as
above. We estimate all parameters by minimising the weighted sum
\begin{equation}\label{eq:joint}
 \mathcal{L}(\theta)=\gamma_{Y}\,\ell_{Y}(\theta)+\gamma_{W}\,\ell_{W}(\theta),
 \qquad \gamma_{Y},\gamma_{W}>0.
\end{equation}
Both terms depend on \(\Psi\) and \(\psi\), so the hidden layer receives gradient contributions from both
edge values and \(H_{ij}\) has to support both heads at once. The single-head neural network has the identical
hidden layer and differs only in that the second term is absent, so comparing the two isolates the cost of
sharing one representation. Figure~\ref{fig:latent} shows the model.


\begin{figure}[htbp]
  \centering
  \resizebox{\textwidth}{!}{%
\begin{tikzpicture}[
      x = 1cm, y = 1cm, font = \small,
      bar/.style  = {draw = black!55, fill = black!6,  line width = .4pt},
      hid/.style  = {draw = black!55, fill = black!22, line width = .4pt},
      hdbox/.style= {draw = black!55, fill = black!38, line width = .4pt},
      arr/.style  = {-{Stealth[length = 2.2mm]}, line width = .7pt},
      grad/.style = {-{Stealth[length = 2mm]}, line width = .6pt, dashed, black!55},
      dim/.style  = {font = \scriptsize, text = black!55, inner sep = 1.5pt},
      note/.style = {font = \scriptsize, inner sep = 2pt}
    ]

  % input ---------------------------------------------------------------
  \draw[bar] (0,-2.10) rectangle (0.55,2.10);
    \node[dim, rotate = 90] at (-0.32,0) {$d$};
    \node[note, below] at (0.27,-2.20) {$z_{ij}$};

  % trunk ---------------------------------------------------------------
  \draw[arr] (0.80,0) -- (3.20,0);
    \node[dim, align = center] at (2.00,0.80)
          {$\Psi\in\mathbb{R}^{r\times d},\ \ \psi\in\mathbb{R}^{r}$};
    \node[dim, align = center] at (2.00,0.40) {$\mathrm{ReLU}$, then dropout};

  % shared representation ------------------------------------------------
  \draw[hid] (3.45,-0.75) rectangle (4.00,0.75);
    \node[dim, rotate = 90] at (3.18,0) {$r$};
    \node[note, below, align = center] at (3.72,-0.85) {$H_{ij}$};

  % two heads ----------------------------------------------------------
  \draw[arr] (4.25,0.32) -- (6.30,1.55);
    \node[dim] at (5.00,1.32) {$\phi_{Y},\,\phi_{0Y}$};
  \draw[arr] (4.25,-0.32) -- (6.30,-1.55);
    \node[dim] at (5.00,-1.32) {$\phi_{W},\,\phi_{0W}$};

  \draw[hdbox] (6.50,1.48) rectangle (6.80,1.78);
  \draw[hdbox] (6.50,-1.78) rectangle (6.80,-1.48);

  \node[note, right] at (7.00,1.63)
        {$\widehat p^{\,Y}_{ij}=\sigma\!\left(\phi_{Y}^{\top}H_{ij}+\phi_{0Y}\right)$};
  \node[note, right] at (7.00,-1.63)
        {$\widehat p^{\,W}_{ij}=\sigma\!\left(\phi_{W}^{\top}H_{ij}+\phi_{0W}\right)$};

  % losses -----------------------------------------------------------------
  \draw[arr] (11.55,1.63) -- (12.25,1.63);
  \draw[arr] (11.55,-1.63) -- (12.25,-1.63);
  \node[note, draw = black!45, fill = black!3, inner sep = 4pt] (ly) at (13.05,1.63)
        {$\ell_{Y}$};
  \node[note, draw = black!45, fill = black!3, inner sep = 4pt] (lw) at (13.05,-1.63)
        {$\ell_{W}$};
  \node[dim, right] at (13.45,1.63) {weighted BCE against $Y_{ij}$};
  \node[dim, right] at (13.45,-1.63) {weighted BCE against $T^{W}_{ij}$};

  \node[note, draw = black!45, fill = black!3, inner sep = 4pt] (loss) at (13.05,0)
        {$\mathcal{L}=\gamma_{Y}\,\ell_{Y}+\gamma_{W}\,\ell_{W}$};
  \draw[grad] (ly.south) -- (loss.north);
  \draw[grad] (lw.north) -- (loss.south);

  % shared gradient ----------------------------------------------------------
  \draw[grad] (loss.west) -- (11.30,0) -- (11.30,-2.90) -- (2.00,-2.90) -- (2.00,-0.20);
    \node[dim, above] at (6.65,-2.87) {both losses update the shared $\Psi,\psi$};

  \end{tikzpicture}%
}
  \caption[The shared-representation model]{The shared-representation model. A single hidden layer maps the
  pair vector \(z_{ij}\in\mathbb{R}^{d}\) to a representation \(H_{ij}\in\mathbb{R}^{r}\) with \(r<d\).
  Two linear heads act on the same \(H_{ij}\) and produce one probability each. Both are scored by a
  weighted binary cross-entropy, and the weighted sum of the two is minimised, so the hidden layer receives
  gradient contributions from both edge values.}
  \label{fig:latent}
\end{figure}
\FloatBarrier

\subsection{Training and Evaluation}
\label{sec:eval}

We partition the pairs \(\mathcal{P}\) once, before fitting any model, into \(F=10\) folds
\(\mathcal{P}_{1},\dots,\mathcal{P}_{F}\) of approximately equal size, stratified on \(Y\) and drawn from
a fixed seed. We fit each model on \(\mathcal{P}\setminus\mathcal{P}_{f}\) and score it on
\(\mathcal{P}_{f}\) for \(f=1,\dots,F\), and every model sees the identical partition. We centre and
scale each coordinate of \(Z\) using the mean and standard deviation of the training fold, except for the
tree ensembles, whose splits are invariant to monotone rescaling of a coordinate.

For the binary edge values we report the area under the receiver operating characteristic curve (AUC) and the
average precision (AP), the latter being the more informative of the two under strong class imbalance
\citep{davis2006}. Both are rank-based, so no decision threshold is applied to \(\widehat p_{ij}\). For
the weighted network we report the coefficient of determination on the original \([0,1]\) scale,
\begin{equation}\label{eq:r2}
 R^{2}_{f}
 =1-\frac{\sum_{(i,j)\in\mathcal{P}_{f}}\bigl(W_{ij}-\widehat W_{ij}\bigr)^{2}}
     {\sum_{(i,j)\in\mathcal{P}_{f}}\bigl(W_{ij}-\bar W_{f}\bigr)^{2}},
 \qquad
 \bar W_{f}=\frac{1}{|\mathcal{P}_{f}|}\sum_{(i,j)\in\mathcal{P}_{f}}W_{ij},
\end{equation}
so that \(R^{2}_{f}<0\) means the fitted values are worse on fold \(f\) than the constant prediction
\(\bar W_{f}\). We report every metric as a mean and a standard deviation over the \(F\) held-out folds. Two models are
compared by the mean of their per-fold differences and, where the consistency of a difference matters, by
the number of folds favouring each model. The training sets of different folds overlap, so a paired test
over the \(F\) differences has no unbiased variance estimator and its \(p\)-values would be
anti-conservative \citep{bengio2004}; we therefore report the differences themselves rather than a test.

\subsection{Representation Comparison}
\label{sec:cka}

Each of the three neural networks produces a hidden layer of width \(r\): write \(H^{Y}_{ij}\) and
\(H^{W}_{ij}\) for the layer of the single-head neural network fitted to \(Y\) and to the binarised edge
weight, and \(H^{YW}_{ij}\) for the shared layer of the joint model. Their coordinates cannot be
compared one by one, since an orthogonal transformation of a hidden layer together with a compensating
change of its head leaves every fitted value unchanged. What can be compared is the geometry each layer
induces on the pairs.

Let \(A\) and \(B\) hold two such layers evaluated on the same \(q\) pairs, stacked column-centred into
\(\mathbb{R}^{q\times r}\). Their linear centred kernel alignment \citep{kornblith2019} is
\begin{equation}\label{eq:cka}
 \operatorname{CKA}(A,B)
 =\frac{\bigl\lVert A^{\top}B\bigr\rVert_{F}^{2}}
       {\bigl\lVert A^{\top}A\bigr\rVert_{F}\,\bigl\lVert B^{\top}B\bigr\rVert_{F}}
 \;\in\;[0,1],
\end{equation}
which equals one when the two layers differ by an orthogonal transformation and a common scaling and is
invariant to both. It is not invariant to rescaling single coordinates, so we scale each coordinate to unit
variance before the comparison; otherwise a hidden unit with large activations would dominate the
alignment for no reason beyond its scale.

Every fold refits each neural network, so layers from different folds are not the same map and we do not
pool them. We compute the three pairwise values within each held-out fold \(\mathcal{P}_{f}\) and report
them as a mean and a standard deviation over the \(F\) folds.

\subsection{Robustness Analyses}
\label{sec:robustness}

We report two further analyses. The first asks whether a gain of the joint model on one network depends on information carried by the other; the second asks whether the
coordinates selected by the penalised linear model are stable under resampling.

\subsubsection{Permutation Control}
\label{sec:shuffley}

To separate a contribution of \(Y\) to the second head from the effect of joint fitting itself, the
we refit the joint model with the \(Y\) labels permuted within each training fold. Let
\(\varrho\) be a uniformly drawn permutation of \(\mathcal{T}\). We then fit the two heads against
\(Y_{\varrho(i,j)}\) and the unpermuted binarised edge weight respectively, with the representation
the hidden layer, the partition, the architecture, the hyperparameters and the random seed unchanged.
Under the permutation, \(Y_{\varrho(i,j)}\) is independent of \(z_{ij}\) by construction, so any
remaining gain on the second head cannot be attributed to information shared between the two
networks.

\subsubsection{Stability Selection}
\label{sec:stability}

To identify coordinates of \(z_{ij}\) that are selected consistently, we apply stability selection
\citep{meinshausen2010} to \(Y\). We resample over nodes rather than over pairs, because two pairs
sharing a node are not independent. For resample \(s\) we draw subsets
\(\mathcal{B}^{(s)}\subset\mathcal{B}\) and \(\mathcal{V}^{(s)}\subset\mathcal{V}\) containing half the
nodes of each set, and the resample consists of the observed pairs they span,
\begin{equation}\label{eq:subsample}
 \mathcal{P}^{(s)}=\bigl\{(i,j)\in\mathcal{P}\;:\;b_{i}\in\mathcal{B}^{(s)},\;v_{j}\in\mathcal{V}^{(s)}\bigr\}.
\end{equation}
Drawing \(K\) complementary pairs \citep{shah2013} gives \(S=2K\) resamples; we redraw a resample with fewer
than five positive \(Y_{ij}\).

Before resampling we remove coordinates that are non-zero in fewer than ten observed pairs or constant over
them, which leaves \(d\) candidate coordinates.

On each resample we centre and scale every coordinate within that resample and fit the
\(L_{1}\)-penalised logistic regression of Equation~\eqref{eq:lasso}, with class weights formed as in
Section~\ref{sec:models}, along the decreasing grid of penalties
\begin{equation}\label{eq:ssgrid}
 \lambda_{t}=\lambda_{\max}\,0.01^{(t-1)/29},\qquad t=1,\ldots,30,
\end{equation}
where \(\lambda_{\max}\) is the smallest penalty at which no coordinate is selected. The path of resample
\(s\) stops at the first penalty \(\lambda^{(s)}\) at which at least \(q\) coordinates are selected, and
\(\widehat S^{(s)}=\{k:\widehat\beta_{k}(\lambda^{(s)})\neq0\}\) records them. The selection probability
of coordinate \(k\) is estimated by the selection frequency
\begin{equation}\label{eq:pisel}
 \widehat\pi_{k}
 =\frac{1}{S}\sum_{s=1}^{S}\mathbf{1}\!\left\{k\in\widehat S^{(s)}\right\},
\end{equation}
and a coordinate is called stable when \(\widehat\pi_{k}\geq\pi_{\mathrm{thr}}\) for a threshold
\(\pi_{\mathrm{thr}}\in(\tfrac{1}{2},1)\). Under the exchangeability assumption of
\citet{meinshausen2010}, the expected number \(V\) of falsely selected coordinates satisfies
\begin{equation}\label{eq:mbbound}
 \mathbb{E}[V]\;\leq\;\frac{\bar q^{\,2}}{\left(2\pi_{\mathrm{thr}}-1\right)d},
\end{equation}
where \(\bar q\) is the average number of coordinates selected per resample. Requiring
\(\mathbb{E}[V]\leq1\) gives the target support size
\begin{equation}\label{eq:qbound}
 q=\left\lfloor\sqrt{\left(2\pi_{\mathrm{thr}}-1\right)d}\right\rfloor,
\end{equation}
at which each path is stopped. Because the grid is discrete, a resample can select slightly more than
\(q\) coordinates, so \(\bar q\) can exceed \(q\). Because the resampling draws nodes rather than
independent observations, we use the bound as a guide to the operating point rather than as an exact
error guarantee.

The penalty and class weighting are those of the penalised linear model of Section~\ref{sec:models}, but
the penalty is set by the stopping rule rather than by cross-validation, so the analysis describes the
coordinates that model class selects.

